%----------------------------------------------------------------------
\section{Assumptions}
%----------------------------------------------------------------------
\textbf{Supplied facts (not assumptions).} 10-minute water levels at S1--S4 and T1; rainfall at
R1 and R2 with timestamps and gauge locations; recorded reservoir releases and announced future
schedules with announcement and effective times; station locations, river connectivity and
along-river distances. A monitored tributary (T1) joins between S1 and S2 and an unmonitored one
between S2 and S3. Stations are unevenly spaced, travel times vary with conditions, precise
travel times and the unmonitored inflow are not supplied, and future rainfall and future levels are
unknown at issuance. Everything below is an assumption made because the brief leaves it open.

\subsection*{Unspecified data conditions}
\begin{description}[leftmargin=0.9cm, style=nextline]
\item[A1 Record length.] At least three years of history are available for all sources and include
several flood seasons. \emph{Why reasonable:} authorities normally keep multi-year telemetry archives.
\emph{Effect on design:} enough events for a small ($\approx$0.3\,M parameter) model, but not for a large
one; data are split chronologically (oldest years train, next season validation, most recent season
test) so that no future information leaks into training.
\item[A2 Reporting delays.] Level telemetry usually arrives within one 10-minute interval but can be up
to 30\,min late; rain gauges report with up to 20\,min delay; reservoir releases are logged within
10\,min. \emph{Why:} typical of radio/GSM telemetry. \emph{Effect:} every input is built ``as of''
$\tau$ with an explicit \emph{age} feature, and training replays these delays (Section~5).
\item[A3 Measurement quality.] Gauges share a common datum; level noise is a few centimetres;
occasional spikes and flat-lines occur. \emph{Why:} standard for pressure/radar gauges.
\emph{Effect:} range and rate-of-change checks flag bad readings, which are then treated as missing.
\item[A4 Missingness.] Gaps are sporadic (a few \% of readings) and occasionally last hours; no
station is missing for most of the record. \emph{Why:} sensors are maintained. \emph{Effect:} masks
plus simple filling suffice; a permanently failed station needs retraining or a fallback model.
\item[A5 Release semantics.] A release value (m$^3$/s) holds from its effective time until the next
record (step function). Announced schedules are usually followed but can be revised or deviate
(gate errors, emergency spills). \emph{Why:} how gate operations are logged. \emph{Effect:} scheduled
and actual releases are separate inputs (Sections~4--5).
\item[A6 Issuance timing.] Forecasts are issued at the grid times
($\ldots$, hh:00, hh:10, $\ldots$) after a short collection cut-off (2\,min), and the target is the
level recorded at $\tau+60$\,min. \emph{Why:} matches the 10-minute logging. \emph{Effect:} defines
$\tau$ and the alignment grid (Section~3).
\end{description}

\subsection*{River dynamics the model assumes}
\begin{description}[leftmargin=0.9cm, style=nextline]
\item[A7 Downstream-only propagation.] There is no tidal influence or significant backwater from
downstream structures on S1--S4. \emph{Why:} an inland river below a reservoir. \emph{Effect:} graph
edges point downstream only, which removes spurious paths and is a strong inductive bias.
\item[A8 Bounded, state-dependent travel times.] The travel time of a flood wave over a distance
$d_{ij}$ is $d_{ij}/c$, with wave celerity $c$ between $c_{\min}=0.5$ and $c_{\max}=4$\,m/s, and $c$
grows with discharge (faster in floods). \emph{Why:} the typical range for lowland and mid-gradient
rivers. \emph{Effect:} edge $j\to i$ may only use lags in
$\Lambda_{ij}=[\lfloor d_{ij}/(c_{\max}\Delta)\rfloor,\ \lceil d_{ij}/(c_{\min}\Delta)\rceil+3]$
(10-minute steps, with a margin). Within $\Lambda_{ij}$ the weights are learned and depend on the
current level (Section~4). Travel time is therefore learned, but physically impossible lags are
excluded.
\item[A9 Smooth, monotone response.] Over one hour, the level change at a station is dominated by
recent upstream changes, routed inflows and local rainfall-runoff. More inflow never lowers the
level downstream. \emph{Why:} mass conservation in an open channel. \emph{Effect:} we predict the
change over persistence (residual form), and an optional penalty keeps the median forecast
non-decreasing in the scheduled release.
\item[A10 Rainfall representativeness.] From the gauge locations, R1 mainly drains to S1, S2 and T1,
and R2 to the unmonitored tributary, S3 and S4. \emph{Why:} nearest-catchment assignment.
\emph{Effect:} initial rainfall-to-node edges. Their weights are learned, so a wrong guess is
down-weighted by the data.
\item[A11 Stationary channel.] No major change of the river bed or the level--flow relation occurs within the
record, and the authority reports such changes. \emph{Why:} short record. \emph{Effect:} one model
for the whole record, with drift monitoring triggering retraining.
\item[A12 Rain after issuance is secondary within one hour.] Rain that has fallen by $\tau$ dominates
the one-hour response, and rain falling after $\tau$ matters mainly for the fastest catchments.
\emph{Why:} runoff concentration times usually exceed one hour. \emph{Effect:} no rainfall forecast
is required, and the remaining effect is absorbed by the upper quantile (Section~6).
\end{description}
